% !TeX root = ../凸优化笔记.tex
% 由主文档通过 \input 载入；请编译主文档。
\setcounter{chapter}{2} % 保留原书章号。
\chapter{数值线性代数与算法实现基础}
\label{app:co-numerical}

优化算法经常把主要成本集中在线性方程求解。数值实现应利用正定性、稀疏性和秩结构，同时区分方程残差与实际解误差。

\section{线性系统、分解与最小二乘}

对可逆方阵，带主元 LU 分解将置换后的矩阵写成下三角与上三角矩阵之积，再作两次三角求解。对称正定矩阵可写成 $LL^{\mathsf T}$，其中 $L$ 为对角元正的下三角矩阵，称为 Cholesky 分解。对称不定矩阵通常用带置换的 $LDL^{\mathsf T}$ 分解，$D$ 可含 $1\times1$ 与 $2\times2$ 对角块。KKT 矩阵即使可逆，也通常不正定，因此不能直接套用普通 Cholesky 分解。

对 $A\in\mathbb R^{r\times n}$ 满列秩的最小二乘，薄 QR 分解为 $A=QR$，$Q^{\mathsf T}Q=I$、$R$ 为可逆上三角矩阵。分解残差可得
\[
 \|Ax-b\|_2^2=\|Rx-Q^{\mathsf T}b\|_2^2+\|(I-QQ^{\mathsf T})b\|_2^2,
\]
所以只需解 $Rx=Q^{\mathsf T}b$。

一般矩阵的奇异值分解为 $A=U\Sigma V^{\mathsf T}$。将非零奇异值取倒数并转置矩形对角结构，得到 $\Sigma^\dagger$，定义 Moore--Penrose 伪逆 $A^\dagger=V\Sigma^\dagger U^{\mathsf T}$。$A^\dagger b$ 是所有最小二乘解中的最小范数解；若方程可行，它就是最小范数精确解。

\begin{example}[秩亏时拟合值可以唯一而参数不唯一]
令 $A=\left[\begin{smallmatrix}1&1\\2&2\end{smallmatrix}\right]$、$b=(1,2)^{\mathsf T}$。所有满足 $x_1+x_2=1$ 的点都使残差为零，因此参数并不唯一，拟合向量 $Ax=b$ 却唯一。把任意解写成 $(1/2,1/2)+(d,-d)$，其平方范数为 $1/2+2d^2$，所以伪逆选择的最小范数解为 $(1/2,1/2)$。不能因为算法返回了这个特定向量，就认为数据唯一确定了两个参数。
\end{example}

\section{块矩阵、Schur 补与 KKT 方程}

对块方程 $Hu+A^{\mathsf T}v=r$、$Au=s$，若 $H$ 可逆，先写 $u=H^{-1}(r-A^{\mathsf T}v)$，再得到 $AH^{-1}A^{\mathsf T}v=AH^{-1}r-s$。式中负号必须与原方程一致；优化中的右端常由负残差组成，机械套公式容易出错。

若 $H\succ0$ 且 $A$ 满行秩，Schur 补正定；若 $H$ 奇异但在 $\ker A$ 上正定，整个 KKT 矩阵仍可能可逆，此时直接的块求逆公式不可用。稀疏直接分解中的填充是原零元素在分解中变为非零的现象，适当排序可降低存储和运算成本。

\section{条件数、舍入误差与规模}

对可逆矩阵 $A$，条件数 $\kappa_2(A)=\|A\|_2\|A^{-1}\|_2$ 衡量输入扰动被放大的程度。若 $Ax=b$、近似解为 $\widehat x$、残差为 $r=b-A\widehat x$，则
\[
 \|x-\widehat x\|_2\leq\|A^{-1}\|_2\|r\|_2,
 \qquad
 \frac{\|x-\widehat x\|_2}{\|x\|_2}
 \leq\kappa_2(A)\frac{\|r\|_2}{\|b\|_2}
\]
在 $x,b\ne0$ 时成立。因此小相对残差未必意味着小相对解误差。后向误差讨论计算解可看成哪个邻近问题的精确解；条件数则决定邻近问题的解与原解能差多少。

例如 $A=\operatorname{diag}(1,10^{-8})$、$b=(1,10^{-8})^{\mathsf T}$ 的真解为 $(1,1)$。近似解 $\widehat x=(1,0)$ 的残差仅为 $(0,10^{-8})^{\mathsf T}$，但第二个分量的误差为一。原因是第二个方向的系数很小，变量的大变化只引起方程左侧的小变化。这与平坦优化目标中小梯度不保证点误差小，是同一种尺度问题的不同表现。

满列秩矩阵满足 $\kappa_2(A^{\mathsf T}A)=\kappa_2(A)^2$。正规方程虽然形式简洁，却可能放大数值困难；QR 和奇异值分解避免显式形成这一平方条件的矩阵。不同方法的实际误差仍与舍入、缩放及算法实现有关。

稠密方阵分解一般需 $O(n^3)$ 运算，已有分解后的每个新右端求解需 $O(n^2)$。稀疏问题取决于非零结构与填充，不能只用原非零元素个数给出统一成本。牛顿法和内点法中若矩阵结构不变，可以复用符号分解；数值因子是否可复用则取决于矩阵值是否改变。

\section{综合练习与结果核验}
\begin{enumerate}
 \item 对 $A=\operatorname{diag}(1,10^{-8})$，构造残差很小但解误差明显的近似解，核对条件数界。
 \item 对同一个满列秩最小二乘问题，比较正规方程、QR 与奇异值分解的结果；逐步使两列接近线性相关，记录残差与参数误差的差别。
 \item 任选一次牛顿或 KKT 步，核验方程残差、方向下降性、更新后可行性及优化停止指标。线性系统求解成功、迭代收敛和达到所需最优性精度是三个需要分别确认的结论。
\end{enumerate}

